% Clase 21 — Belmont Muñoz Samuel (botones y enlace USB HID)
\paragraph{Botones y enlace USB HID}
\textbf{Autor:} Belmont Muñoz Samuel

\textbf{Botones:} driver propio en la Pico (\texttt{read\_74hc165.c} y \texttt{read\_74hc165.pio}). El genérico no alcanza:
el periférico SPI no genera el pulso de carga paralela (LD) que el 74HC165 necesita antes de cada lectura. La PIO genera
LD y el reloj, arma tramas de 10 bits con \textit{autopush}, y dos canales DMA encadenados
(\texttt{read\_74hc165.c}, línea 60) las copian a \texttt{g\_controls.buttons} sin CPU. Cómo compilarlo y grabarlo
está en \texttt{firmware/README.md}.

\textbf{Enlace USB:} en la Pico, la pila TinyUSB en modo dispositivo HID (\texttt{usb\_data\_tx.c}). En la Pi no hace
falta driver propio: el kernel reconoce el descriptor HID y carga los genéricos \texttt{usbhid} y \texttt{hid-generic}, que
publican el control en \texttt{/dev/input/eventX} (lo lee RetroArch) y en \texttt{/dev/hidrawX}
(por ahí llegan las órdenes de actuadores).

\textbf{Verificación} (con el control conectado):
\begin{lstlisting}
lsusb | grep -i cafe
cat /proc/bus/input/devices | grep -A4 Arcade
\end{lstlisting}
